\scanpage{15}
\mspage{7}
voir, on note que \inserted{dans} la description (19) d'un
$\widehat C_0$-morphisme (18), $Y'$ est un revêtement (topologique)
étale fini, qui provient canoniquement
d'un rev. étale fini schématique; quant à
$\widehat f_0$, on peut l'interpréter comme un
``isomorphisme'' \deleted{\unsure{de}} \inserted{ou ``classe de chemins''} entre des points (à coefficients
dans $\mathbb C$) d'une multiplicité modulaire convenab-
le, au sens du groupoïde fondamental
de celle-ci. De $=$, le diagramme (21),
montre que la composition des morphismes
dans $\widehat C_0$ s'explicite en termes de composition
de morphismes, d'une part dans le groupoïde
fondamental de la multiplicité modulaire
pour les courbes anabéliennes, d'autre part
par des morphismes étales finis entre courbes
algébriques. La construction de la catégorie
``isotopique profinie'' ad-hoc des courbes algébriques anabéliennes
sur un corps alg. clos $k$, peut d'ailleurs
se faire sur tt corps alg. clos $k$
(pas $=$ besoin qu'il soit de car. $0$ !), et
cette catégorie dépend fonctoriellement de $k$,
\margin{14}
